\section{Modifications to clause 9 (Declarations)}

The following modifications shall be applied to clause 9 (\emph{Declarations}
[\texttt{dcl}]):

\begin{modification}[{Removal of 9.2.9.7.1[dcl.spec.auto.general]/14}]
\stdremove{A result binding never has an undeduced placeholder type (9.4.2).}
\newline

\textbf{Note to Editors}: Also \textbf{remove completely} example 5.
\end{modification}

\begin{modification}[{Modification of 9.3.1[dcl.decl.general]/1}]
 A declarator declares a single variable, function, or type, within a
declaration. The \emph{init-declarator-list} appearing in a
\emph{simple-declaration} is a \emph{comma-separated} sequence of declarators,
each of which can have an initializer.\\

init-declarator-list :\\
\tabto{1em}init-declarator\\
\tabto{1em}init-declarator-list , init-declarator\\

init-declarator :\\
\tabto{1em}declarator initializer\\
\tabto{1em}declarator requires-clause\textsubscript{opt}
\stdremove{function-contract-specifier-seq\textsubscript{opt}}
\end{modification}

\begin{modification}[{Removal of 9.3.1[dcl.decl.general]/5}]
\stdremove{The optional \emph{function-contract-specifier-seq} (9.4.1) in an
\emph{init-declarator} shall be present only if the declarator
declares a function.}
\end{modification}

\begin{modification}[{Removal of section 9.4[dcl.contract]}]
Section 9.4 (\emph{Function contract specifiers} [\texttt{dcl.contract}]) shall be 
\textbf{completely removed}.
\end{modification}

\begin{modification}[{Modification of 9.6.1[dcl.fct.def.general]/1}]
Function definitions have the form\\

function-definition :\\
\tabto{1em}attribute-specifier-seq\textsubscript{opt} decl-specifier-seq\textsubscript{opt} declarator
virt-specifier-seq\textsubscript{opt}\\
\tabto{2em}\stdremove{function-contract-specifier-seq\textsubscript{opt}} function-body\\
\tabto{1em}attribute-specifier-seq\textsubscript{opt} decl-specifier-seq\textsubscript{opt} declarator
requires-clause\\
\tabto{2em}\stdremove{function-contract-specifier-seq\textsubscript{opt}} function-body\\

function-body :\\
\tabto{1em}ctor-initializer\textsubscript{opt} compound-statement\\
\tabto{1em}function-try-block\\
\tabto{1em}= \texttt{default} ;\\
\tabto{1em}deleted-function-body\\

deleted-function-body :\\
\tabto{1em}= \texttt{delete} ;\\
\tabto{1em}= \texttt{delete} ( unevaluated-string ) ;\\
\end{modification}

\begin{modification}[{Removal of 9.6.4[dcl.fct.def.coroutine]/6 Note 1}]
\stdremove{NOTE 1 An odr-use of a non-reference parameter in a postcondition assertion of a
coroutine is ill-formed (9.4.1).}
\end{modification}

\begin{modification}[{Modification of 9.13.1[dcl.attr.grammar]}]
Attributes and annotations specify additional information for various source constructs such as types, variables,
names, \stdremove{contract assertions,} blocks, or translation units.\\

\ldots
\end{modification}

\begin{modification}[{Modification of 9.13.8[dcl.attr.unused]/2}]
The attribute may be applied to the declaration of a class, type alias, 
variable (including a structured binding
declaration), structured binding, \stdremove{result binding (9.4.2),} 
non-static data member, function, enumeration, or
enumerator, or to an identifier label (8.2).
\end{modification}
